\section{USAMO 1980}
        \begin{enumerate}
        	\item (\textit{USAMO 1980}) A balance has unequal arms and pans of unequal weight. It is used to weigh three objects. The first object balances against a weight $A$, when placed in the left pan and against a weight $a$, when placed in the right pan. The corresponding weights for the second object are $B$ and $b$. The third object balances against a weight $C$, when placed in the left pan. What is its true weight?


 



\textbf{Solution}

The effect of the unequal arms and pans is that if an object of weight $x$ in the left pan balances an object of weight $y$ in the right pan, then $x = hy + k$ for some constants $h$ and $k$. Thus if the first object has true weight x, then $x = hA + k, a = hx +  k$. 

 So $a = h^2A + (h+1)k$. 

 Similarly, $b = h^2B + (h+1)k$. Subtracting gives $h^2 = \frac{a-b}{A-B}$ and so

 
 \[(h+1)k = a - h^2A = \frac{bA - aB}{A - B}\].

 The true weight of the third object is thus:

 
 \[hC + k = \\ \boxed{\sqrt{ \frac{a-b}{A-B}} C + \frac{bA - aB}{(A - B)(\sqrt{ \frac {a-b}{A-B}}+ 1)}}\].

 More readably: 
 \[\boxed{ h=\sqrt{\frac{a-b}{A-B}} ;  \\ \text{weight} = hC + \frac{bA - aB}{(A - B)(h + 1)}}\]

 Credit: John Scholes \href{https://prase.cz/kalva/usa/usoln/usol801.html}{https://prase.cz/kalva/usa/usoln/usol801.html}

 


	\item (\textit{USAMO 1980}) Find the maximum possible number of three term arithmetic progressions in a monotone sequence of $n$ distinct reals.


 



\textbf{Solution}

Consider the first few cases for $n$ with the entire $n$ numbers forming an arithmetic sequence 
 \[(1, 2, 3, \ldots, n)\]If $n = 3$, there will be one ascending triplet (123). Let's only consider the ascending order for now.If $n = 4$, the first 3 numbers give 1 triplet, the addition of the 4 gives one more, for 2 in total.If $n = 5$, the first 4 numbers give 2 triplets, and the 5th number gives 2 more triplets (135 and 345).Repeating a few more times, we can quickly see that if $n$ is even, the nth number will give 
 \[\frac{n}{2} - 1\] more triplets in addition to all the prior triplets from the first $n-1$ numbers. If $n$ is odd, the $n$th number will give 
 \[\frac{n-1}{2}\] more triplets. Let $f(n)$ denote the total number of triplets for $n$ numbers. The above two statements are summarized as follows: If $n$ is even, 
 \[f(n) = f(n-1) + \frac{n}2 - 1\]If $n$ is odd, 
 \[f(n) = f(n-1) + \frac{n-1}2\] 

 Let's obtain the closed form for when $n$ is even:
 \begin{align*} f(n) &= f(n-2) + n-2\\ f(n) &= f(n-4) + (n-2) + (n-4)\\ f(n) &= \sum_{i=1}^{n/2} n - 2i\\ \Aboxed{f(n\ \text{even}) &= \frac{n^2 - 2n}4} \end{align*}


 Now obtain the closed form when $n$ is odd by using the previous result for when $n$ is even:
 \begin{align*} f(n) &= f(n-1) + \frac{n-1}2\\ f(n) &= \frac{{(n-1)}^2 - 2(n-1)}4 + \frac{n-1}2\\ \Aboxed{f(n\ \text{odd}) &= \frac{{(n-1)}^2}4} \end{align*}


 Note the ambiguous wording in the question! If the "arithmetic progression" is allowed to be a disordered subsequence, then every progression counts twice, both as an ascending progression and as a descending progression.Double the expression to account for the descending versions of each triple, to obtain:
 \begin{align*} f(n\ \text{even}) &= \frac{n^2 - 2n}2\\ f(n\ \text{odd}) &= \frac{{(n-1)}^2}2\\ \Aboxed{f(n) &= \biggl\lfloor\frac{(n-1)^2}2\biggr\rfloor} \end{align*}


 



\textbf{Solution 2}

Unlike in the previous solution, I am interpreting "arithmetic progressions" to be in the same direction as the monotone sequence of real numbers. So, for example, if the sequence of reals is $1 < 4 < 5 < 6$, then $(4, 5, 6)$ counts but $(6, 5, 4)$ doesn't.

 The first solution proves that there is a sequence with $\lfloor (n - 1)^2 / 4 \rfloor$ such arithmetic progressions. In this solution, I prove that this is in fact the upper bound.

 Count the arithmetic progressions by their center values. Let $x_k$ be the $k$th greatest number. The number $x_k$ is the center of at most $\min(k - 1, n - k)$ arithmetic progressions, because any such progression includes a number less than $x_k$ and a number greater than $x_k$. There are $k - 1$ numbers less than $x_k$ and $n - k$ numbers greater than $x_k$, and none of these can be in two separate arithmetic progressions with the same middle term $x_k$.

 Thus, the total number of arithmetic progressions is at most
 \[\sum_{k = 1}^n \min(k - 1, n - k).\]This equals
 \[1 + 2 + \cdots + \frac{n - 3}{2} + \frac{n - 1}{2} + \frac{n - 3}{2} + \cdots + 2 + 1 = \frac{(n - 1)^2}{4}\]if $n$ is odd, and
 \[1 + 2 + \cdots + \frac{n-2}{2} + \frac n2 + \frac n2 + \frac{n-2}{2} + \cdots + 2 + 1 = \frac{(n - 1)^2 - 1}{4}\]if $n$ is even. This can be summarized as
 \[\left\lfloor \frac{(n - 1)^2}{4} \right\rfloor.\]

 


	\item (\textit{USAMO 1980}) $A + B + C$ is an integral multiple of $\pi$. $x, y,$ and $z$ are real numbers. If $x\sin(A)+y\sin(B)+z\sin(C)=x^2\sin(2A)+y^2\sin(2B)+z^2\sin(2C)=0$, show that $x^n\sin(nA)+y^n \sin(nB) +z^n \sin(nC)=0$ for any positive integer $n$.


 



\textbf{Solution}

Let $a=xe^{iA}$, $b=ye^{iB}$, $c=ze^{iC}$ be numbers in the complex plane. 

 Note that $A+B+C=k\pi$ implies $abc=xyz(e^{ik\pi})=\pm xyz$ which is real. Also note that $x\sin(A), y\sin(B), z\sin(C)$ are the imaginary parts of $a, b, c$ and that $x^2\sin(2A), y^2\sin(2B), z^2\sin(2C)$ are the imaginary parts of $a^2, b^2, c^2$ by de Moivre's Theorem. Therefore, $a+b+c$ and $a^2+b^2+c^2$ are real because their imaginary parts sum to zero.

 Finally, note that $\frac{1}{2}\left((a+b+c)^2-(a^2+b^2+c^2)\right)=ab+bc+ac$ is real as well.

 It suffices to show that $P_n=a^n+b^n+c^n$ is real for all positive integer $n$, which can be shown by induction.

 Newton Sums gives the following relationship between sums of the form $P_k=a^k+b^k+c^k$
 \[P_k-S_1P_{k-1}+S_2P_{k-2}-S_3P_{k-3}=0\]Where $S_1=a+b+c$, $S_2=ab+bc+ac$, and $S_3=abc$. It is given that $P_0, P_1, P_2$ are real. Note that if  $P_{k-1}, P_{k-2}, P_{k-3}$ are real, then clearly $P_k$ is real because all other parts of the above equation are real, completing the induction.

 


	\item (\textit{USAMO 1980}) The insphere of a tetrahedron touches each face at its centroid. Show that the tetrahedron is regular.


 



\textbf{Solution}

Let $ABCD$ be the tetrahedron, and let $P$ and $Q$ be the points at which the insphere touches faces $ABC$ and $ABD$ respectively (and therefore the centroids of those faces). Looking at the plane containing $P$, $A$, and $Q$, we see that the intersection of the sphere and the plane is a circle, and that $AP$ and $AQ$ are both tangent to said circle. $[AP]$ and $[AQ]$ are tangents to the circle from the same point and thus have the same length. The same goes for $[BP]$ and $[BQ]$. Thus, triangles $ABP$ and $ABQ$ are congruent, and $\mathbf{\angle PAB = \angle QAB}$.

 Let the intersection of $AP$ and $BC$ be $M$, and let the intersection of $AQ$ and $BD$ be $N$. Then $AM$ and $AN$ are medians of $ABC$ and $ABD$, and thus $|AM| = \frac{3}{2}|AP| = \frac{3}{2}|AQ| = |AN|$. We already know from the previous congruence that $\angle{MAB} = \angle{PAB} = \angle{QAB} = \angle{NAB}$, and $|AB|$ is equal to itself. Thus, $MAB$ and $NAB$ are also congruent to each other. Finally, $|BC| = 2|BM| = 2|BN| = |BD|$ (Because $M$ and $N$ are midpoints of $BC$ and $BD$ respectively), and from the congruence of $MAB$ and $NAB$ we have $\angle{CBA} = \angle{MBA} = \angle{NBA} = \angle{DBA}$, and again $|AB|$ is equal to itself. Thus $ABC$ and $ABD$ are congruent, thus $\mathbf{|AC| = |AD|}$ and $\mathbf{|BC| = |BD|}$.

 By applying the same logic to faces $BCD$ and $ACD$ we get $|AC| = |BC|$ and $|AD| = |BD|$. Finally, applying the same logic to faces $ACB$ and $ACD$ we get $|AB| = |AD|$ and $|CB| = |CD|$. Putting all these equalities together, we get that all edges of the tetrahedron are equal, and thus the tetrahedron is regular. $\blacksquare$

 


	\item (\textit{USAMO 1980}) If $x, y, z$ are reals such that $0\le x, y, z \le 1$, show that $\frac{x}{y + z + 1} + \frac{y}{z + x + 1} + \frac{z}{x + y +  1} \le 1 - (1 - x)(1 - y)(1 - z)$


 



\textbf{Solution}

Rewrite the given inequality so that $1$ is isolated on the right side. Set the left side to be $f(x, y, z)$. Now a routine computation shows

 
$\frac{\partial^2 f}{\partial x^2} = \frac{2y}{(x + z + 1)^3} + \frac{2z}{(x + y + 1)^3}\geq 0$

 
which shows that $f$ is convex (concave up) in all three variables. Thus the maxima can only occur at the endpoints, i.e. if and only if $x, y, z \in \{0,1\}$. Checking all eight cases shows that the value of the expression cannot exceed 1.

 


\end{enumerate}